\newpage
\section{BENCHMARK}
Através da pesquisa de soluções já existentes no mercado, descobrimos que as soluções atuais focam em acessibilidade visual (botões grandes, interfaces limpas) e em conectar a pessoa idosa à sua família ou a cuidadores de forma individualizada. A maioria das plataformas analisadas falha no aspecto comunitário e de gestão logística. O nosso cenário descreve um sistema que empodera mediadores/tutores comunitários a gerenciarem grupos, criarem atividades sociais e acompanhar relatórios de parcipação. Nenhuma das plataformas pesquisadas possui a combinação de:

\begin{enumerate}
    \item Painel administrativo para criação e gestão de inscrições em atividades.
    \item Ambiente de grupo seguro e mediado.
    \item Videochamadas nativas integradas à agenda da comunidade.
\end{enumerate}

Apresentamos um resumo das plataformas pesquisadas nos próximos tópicos.

\subsection{Projeto Rede Bem Estar (App Madu - Instituto Tellus)}
Esta é uma plataforma digital (site e aplicativo) focada em entregar conteúdos relevantes sobre o envelhecimento para idosos, familiares e cuidadores, com o objetivo de fortalecer vínculos e quebrar estereótipos.

\begin{itemize}
    \item \textbf{Requisitos Funcionais Contemplados}
        \begin{itemize}
            \item Acesso a conteúdos compartilhados: O app possui categorias de notícias, dicas de saúde, finanças e vídeos.
            \item Lembretes automáticos: Possui um sistema de notificações a cada três dias baseado nos interesses do usuário.
            \item Cadastro de usuários: Permite criação de perfil para personalização de conteúdo.
        \end{itemize}
    \item \textbf{Requisitos Funcionais NÃO Contemplados}
        \begin{itemize}
            \item Não possui visualização ou inscrição em atividades ao vivo.
            \item Não possui ferramenta nativa para videochamadas.
            \item Não permite troca de mensagens em grupo dentro do app (no projeto piloto, utilizaram o WhatsApp à parte).
            \item Não possui painel de gestão para mediadores criarem/editarem atividades ou gerenciarem listas de participantes.
        \end{itemize}
    \item \textbf{Requisitos Não-Funcionais Contemplados}
        \begin{itemize}
            \item Interface simples e intuitiva: Foco em baixa complexidade e facilidade de leitura.
            \item Funciona em celulares, tablets e computadores: Possui versão web, iOS e Android.
            \item Compatibilidade de acessibilidade: O site conta com o widget VLibras.
        \end{itemize}
    \item \textbf{Diferenciais / Funcionalidades Extras}
        \begin{itemize}
            \item \textbf{Identidade visual humanizada:} Uso da persona "Madu" para gerar conexão com o usuário.
            \item \textbf{Kit do Gestor:} Um programa educativo e materiais de formação (online e impresso) voltado para capacitar facilitadores e gestores de grupos de idosos, o que dialoga muito com a sua ideia de "mediadores comunitários".
        \end{itemize}
\end{itemize}

\subsection{Memoryboard}
Trata-se de um quadro de mensagens digital em formato de hardware (um display de 10 ou 15 polegadas) que fica na casa do idoso. A família gerencia o que aparece na tela por meio de um aplicativo remoto.

\begin{itemize}
    \item \textbf{Requisitos Funcionais Contemplados}
        \begin{itemize}
            \item Acesso a conteúdos compartilhados: Recebe fotos e mensagens.
            \item Lembretes automáticos: A família pode programar lembretes (ex: "hora do remédio").
            \item Troca de mensagens: Recebe mensagens de familiares.
        \end{itemize}
    \item \textbf{Requisitos Funcionais NÃO Contemplados}
        \begin{itemize}
            \item A comunicação é praticamente unilateral (da família para o idoso). Não há inscrição em atividades comunitárias, chat entre vários idosos, videochamadas integradas ou gestão de grupos.
        \end{itemize}
    \item \textbf{Requisitos Não-Funcionais Contemplados}
        \begin{itemize}
            \item Interface simples e intuitiva: Extrema facilidade de uso, pois o idoso não precisa interagir tecnicamente com o aparelho.
            \item Textos legíveis: Design focado em exibição clara e direta.
        \end{itemize}
    \item \textbf{Diferenciais / Funcionalidades Extras}
        \begin{itemize}
            \item \textbf{Remoção total da barreira tecnológica:} Como é um hardware dedicado de exibição passiva, resolve o problema de idosos com severa dificuldade motora ou cognitiva que não conseguem operar um smartphone comum.
        \end{itemize}
\end{itemize}

\subsection{Oscar Senior}
Um aplicativo de comunicação projetado especificamente para adultos mais velhos, focado em segurança e chamadas simplificadas.

\begin{itemize}
    \item \textbf{Requisitos Funcionais Contemplados}
        \begin{itemize}
            \item Acesso a atividades online: Possui videochamadas nativas e intuitivas.
            \item Troca de mensagens: Permite mensagens instantâneas.
            \item Lembretes automáticos: Focado em lembretes de medicação e rotina.
        \end{itemize}
    \item \textbf{Requisitos Funcionais NÃO Contemplados}
        \begin{itemize}
            \item Não é uma plataforma de "comunidade". Faltam recursos para mediadores criarem eventos, gerenciarem inscrições e compartilharem links de atividades em grupo.
        \end{itemize}
    \item \textbf{Requisitos Não-Funcionais Contemplados}
        \begin{itemize}
            \item Interface simples e intuitiva: Botões grandes e navegação facilitada.
            \item Autenticação segura: Cria um ambiente seguro e fechado para a família.
        \end{itemize}
    \item \textbf{Diferenciais / Funcionalidades Extras}
        \begin{itemize}
            \item \textbf{Alertas de emergência (SOS):} Um botão de pânico integrado, caso o idoso precise de ajuda imediata.
        \end{itemize}
\end{itemize}

\subsection{Connect Care App}
Aplicativo focado na conexão entre o idoso, seus familiares e profissionais de saúde/cuidadores.

\begin{itemize}
    \item \textbf{Requisitos Funcionais Contemplados}
        \begin{itemize}
            \item Videochamadas e Troca de mensagens seguras.
            \item Lembretes automáticos: Notificações de saúde.
        \end{itemize}
    \item \textbf{Requisitos Funcionais NÃO Contemplados}
        \begin{itemize}
            \item Assim como o Oscar Senior, falta o módulo de "Gestão de Atividades" (criação de turmas, relatórios de participação, inscrição em eventos comunitários).
        \end{itemize}
    \item \textbf{Diferenciais / Funcionalidades Extras}
        \begin{itemize}
            \item \textbf{Agendamento de compromissos e Check-ins de bem-estar:} Funcionalidades excelentes para tutores monitorarem a saúde e o humor diário dos idosos participantes.
        \end{itemize}
\end{itemize}

\subsection{Silver Surf Web Browser}
Não é uma plataforma social, mas sim um navegador de internet adaptado para a terceira idade.

\begin{itemize}
    \item \textbf{Como atende aos requisitos:} Ele brilha nos \textbf{Requisitos Não Funcionais.} O Silver Surf atende aos requisitos de Botões grandes, textos legíveis, possibilidade de aumentar a fonte, suporte a alto contraste e navegação minimizada.
    \item \textbf{Diferenciais / Funcionalidades Extras}
        \begin{itemize}
            \item Permite que o idoso acesse outras redes sociais (como Facebook ou portais de notícias) com uma interface sobreposta que melhora a acessibilidade geral da internet. É uma excelente inspiração de UI/UX (Design de Interface) para o seu projeto.
        \end{itemize}
\end{itemize}